\subsection{The cyclic quotient at thirty-seven and its remaining arithmetic question}
\label{subsec:thirtyseven-gate}

The order-three symmetry of a cyclic Fermat quotient at thirty-seven
produces a rational sixfold, but that sixfold has positive rank.
The complementary twelvefold is the remaining rank-zero candidate
within this particular cyclic quotient.  The following calculation
does not establish a new section classification.

Let $C/\mathbb Q$ be the smooth proper model of
\[
 y^{37}=x(1-x)^{10},\qquad J=\operatorname{Jac}(C),
\]
and put
\[
 c(x,y)=\left(\frac1{1-x},\frac{y^{10}}{(1-x)^3}\right),\quad
 B=C/\langle c\rangle,\quad A=\operatorname{Jac}(B),\quad
 P=\operatorname{im}(2-c-c^2)\subset J.
\]
Let $R_0,R_1,R_\infty$ be the unique points above $0,1,\infty$.
Write
\[
 K=\mathbb Q(\zeta_{37}),\quad H=\langle10\rangle=\{1,10,26\},
 \quad M=K^H,\quad F_0=M^+.
\]

\begin{proposition}[The rational factors at thirty-seven]
\label{prop:thirtyseven-factors}
There is a rational isogeny $J\sim_{\mathbb Q}A\times P$, with
$\dim A=6$, $\dim P=12$, and
\[
 \operatorname{End}_{\mathbb Q}^0(J)
 =\mathbb Q\times\mathbb Q(\zeta_3),\qquad
 \operatorname{End}_{\mathbb Q}^0(A)=\mathbb Q,\qquad
 \operatorname{End}_{\mathbb Q}^0(P)=\mathbb Q(\zeta_3).
\]
Both factors are rationally simple, and
\[
 \operatorname{rank}A(\mathbb Q)\geq1,\qquad
 \operatorname{rank}P(\mathbb Q)\in2\mathbb Z_{\geq0},\qquad
 P(\mathbb Q)_{\mathrm{tors}}\simeq\mathbb Z/37\mathbb Z.
\]
Thus $P$ is the only possible positive-dimensional rational rank-zero
isogeny factor of this $J$.
\end{proposition}
\begin{proof}
The degree-thirty-seven map to the $x$-line is fully ramified at
$0,1,\infty$, so $g(C)=18$.  The identity
$10^2+10+1=3\cdot37$ gives, by direct substitution,
\[
 c^2(x,y)=\left(\frac{x-1}{x},
       -\frac{y^{26}}{x(1-x)^7}\right),\qquad c^3=1.
\]
The three branch points are cyclically permuted.  The other fixed
points have $x^2-x+1=0$.  For $\rho+\rho^{-1}=1$, the curve and
fixed-point equations become $y^{37}=y^9=-1$, whose unique common
solution is $y=-1$.  Thus there are exactly two fixed points,
$(\rho,-1)$ and $(\rho^{-1},-1)$, and
\[
 34=3(2g(B)-2)+4.
\]
This proves $g(B)=6$.  Norm and pullback identify the image of
$(1+c+c^2)/3$ with $A$ up to isogeny and give the stated complement.

The Jacobi-sum CM type on $K$ is
\[
 S=\{1,2,3,4,5,6,8,9,10,12,13,15,16,19,20,23,26,30\};
\]
equivalently, $b\in S$ if $b+(10b\bmod37)<37$.
Its multiplicative stabilizer is exactly $H$.  With generator
$2\bmod37$, the descended type on the cyclic degree-twelve field
$M$ is
\[
 T=\{0,1,2,3,4,11\}\subset\mathbb Z/12\mathbb Z,
\]
whose translation stabilizer is trivial: a preserving translation
must belong to $T$, and each of $1,2,3,4,11$ sends an element of
$T$ outside $T$. The CM decomposition of
\citet[Lemma~1.6 and the following remark, pp.~206--207]{GrossRohrlich1978}
therefore gives $J_{\overline{\mathbb Q}}\sim D^3$, with $D$
absolutely simple of dimension six and
$\operatorname{End}_{\overline{\mathbb Q}}^0(D)=M$.
If $\tau(x,y)=(x,\zeta_{37}y)$, then $c\tau c^{-1}=\tau^{10}$.
The crossed product generated by $K$ and $c$ is
\[
 K\rtimes H=\operatorname{End}_M(K)=M_3(M)
          =\operatorname{End}_{\overline{\mathbb Q}}^0(J).
\]
Galois acts on the unique expansion $\sum_{j=0}^2a_jc^j$ through
its usual action on the $a_j\in K$ and fixes $c$.  The rational
fixed algebra is $\mathbb Q[c]=\mathbb Q\times\mathbb Q(\zeta_3)$.
Its central factors prove the endomorphism assertions and rational
simplicity.  The same argument over $F_0$ gives $F_0[c]$, so
$A_{F_0}$ has real multiplication by $F_0$.

On $X^{37}+Y^{37}=1$, the quotient map is
$(X,Y)\mapsto(X^{37},XY^{10})$.  It takes
$(\rho,\rho^{-1})$ to $(\rho,-1)$.
\citet[Theorem~2.1, pp.~208--209]{GrossRohrlich1978} applies because
$10\notin\{1,18,35\}$, and proves that
\[
 D_0=[(\rho,-1)+(\rho^{-1},-1)-2R_\infty]\in J(\mathbb Q)
\]
has infinite order.  Changing the normalizing branch cusp changes
this class only by $37$-torsion.  The two displayed points are fixed
by $c$, whereas differences of branch points are killed by $37$.
Consequently, for $\pi:C\to B$,
\[
 \pi^*\pi_*D_0=(1+c+c^2)D_0=3D_0+\text{torsion}
\]
is non-torsion.  Thus $A(\mathbb Q)$ has positive rank.
The $\mathbb Q(\zeta_3)$-action makes the rank of $P(\mathbb Q)$ even.

By \citet[Theorem~1.1, p.~202]{GrossRohrlich1978},
$J(\mathbb Q)_{\mathrm{tors}}\simeq C_{37}$.
Put $D=[R_0-R_\infty]$ and $E=[R_1-R_\infty]$.
The divisor of $x$ gives $37D=0$, and $D\ne0$, since a principal
divisor $R_0-R_\infty$ would give a degree-one map to the projective
line.  The divisor of $y$ gives $D+10E=0$, so
$cD=E-D=10D$ and $c^2D=26D$.
Hence $(2-c-c^2)D=3D$, which generates the order-thirty-seven
torsion and belongs to $P$.  The inclusion $P\subset J$ proves
the torsion assertion.
\end{proof}

For an auxiliary prime $\ell\ne37$ and an abelian variety $B'$,
write $V(B')=T_\ell B'\otimes_{\mathbb Z_\ell}\mathbb Q_\ell$.
Write $N$ for the arithmetic conductor and $W$ for the global root number.

\begin{proposition}[Conductor, signs, and the order-nine character]
\label{prop:thirtyseven-character}
The conductors and root numbers are
\[
 N(J)=37^{36},\quad N(A)=37^{12},\quad N(P)=37^{24},
 \qquad W(J)=W(A)=-1,\quad W(P)=+1.
\]
Let $\chi$ be the character modulo $37$ with $\chi(2)=\zeta_9$.
Over a coefficient field containing its values,
\[
 V(J)\simeq V(A)\oplus V(A)\otimes\chi
                   \oplus V(A)\otimes\chi^2.
\]
The two complementary weight-two Hilbert constituents over the real
sextic field $F_0$ have nontrivial cubic central characters.
\end{proposition}
\begin{proof}
In the parameters $(r,s,t,\delta)=(1,10,26,1)$,
\[
 10^{10}(-11)^{26}\equiv581\pmod{37^2},\qquad
 581^{36}\equiv1\pmod{37^2}.
\]
Thus \citet[Proposition~3.12, p.~16; Theorem~3.13,
pp.~17--18]{ShuExpanded2021} applies in the tame case.
The Jacobi-sum character on $K$ has conductor the prime above $37$
and nontrivial quadratic tame inertia.  Its induction to $\mathbb Q$
is tame and has no inertia invariants, so its conductor exponent is
$36$.  Each rational direct summand is also tame and has no inertia
invariants.  Its exponent is consequently its representation
dimension, giving $12$ and $24$.  The same root-number formula gives
$W(J)=(2/37)=-1$.

The equalities $26=2^{12}$ and $10=2^{24}\pmod{37}$ show that every
cubic character modulo $37$ is trivial on $H$.  The chosen
order-nine $\chi$ instead restricts nontrivially to $H$.
Let $\varphi$ be the CM Hecke character of $M$ associated with $A$
in one fixed embedding, and put $\nu=\chi\circ N_{M/\mathbb Q}$.
The latter cuts out $K/M$.  The $c$-invariant line extends the
Jacobi-sum character from $G_K$ to $G_M$; the three extensions are
$\varphi,\varphi\nu,\varphi\nu^2$.  Transitivity of induction gives
the asserted representation identity and
\[
 L(P/\mathbb Q,s)=L_M(\varphi\nu,s)L_M(\varphi\nu^2,s),
 \qquad L(A/\mathbb Q,s)=L_M(\varphi,s).
\]
Since $\chi^3|_{G_M}=1$, it is a self-twist of $V(A)$.
Thus $V(A)\otimes\chi^{-1}\simeq V(A)\otimes\chi^2$.
Polarization makes the two normalized complementary constituents a
dual pair.  Their completed functional equations have reciprocal
root numbers, so $W(P)=+1$ and $W(A)=W(J)=-1$.

Quadratic induction from $M/F_0$ gives a weight-two Hilbert form
$\pi$ with trivial normalized central character.  Indeed the real
multiplication on $A_{F_0}$ is fixed by Rosati, so every
two-dimensional constituent has cyclotomic determinant.
The character $\theta=\chi\circ N_{F_0/\mathbb Q}$ has order three,
since $G_{F_0}$ has image $\langle2^6\rangle$ in
$\operatorname{Gal}(K/\mathbb Q)$.  The complementary forms
$\pi\otimes\theta$ and $\pi\otimes\theta^2$ therefore have central
characters $\theta^2$ and $\theta$, respectively.
\end{proof}

The arithmetic conclusion needed for a rank-zero finite-descent
application remains unproved.  In particular, no assertion is made
that $L(P/\mathbb Q,1)\ne0$, that $\operatorname{rank}P(\mathbb Q)=0$,
or that $\Sha(P/\mathbb Q)$ is finite.  The following applicability
restrictions explain why the previous arithmetic inputs do not
immediately settle it.

The direct central-value statements of
\citet[Theorem~A, p.~28]{Zhang2001} and
\citet[Theorems~B and~B', pp.~593--594]{Nekovar2012} require trivial
central character.  The two complementary forms violate this
condition.  Nekov\'a\v r's more general Theorem~A permits a
compensating character over an auxiliary CM extension, subject to
equation~(0.3.1), but this changes the Selmer representation.
Taking that character trivial again requires trivial central
character.  Killing the cubic twists by passing to $K^+$ makes
$P$ isogenous to $A^2$ and imports its positive-rank point.

The quadratic-Hilbert construction of \citet{LoefflerZerbes2025},
including the rational refinement used above, does not give a theorem
for this real sextic base.  The primitive type cannot be reduced to
the quartic CM subfield: its counts above the four embeddings are
$(2,1,1,2)$, whereas an induced type would have counts zero or three.
There is no imaginary quadratic subfield of the cyclic degree-twelve
field $M$; its unique quadratic subfield is real.

In the September 2025 author version of \citet{Hsieh2014},
Theorem~2 and Corollaries~2--3, pp.~3--4, the general inclusion
$(F_{\psi,\Sigma})\subset(L_{\psi,\Sigma})$ has the lower-bound
direction for the Selmer module.  The anticyclotomic equality case
fails here at $37$: complex conjugation fixes $\nu$, and
\[
 (\varphi\nu)^c(\varphi\nu)=\operatorname{Norm}\,\nu^2.
\]
The nontrivial cubic tame inertia of $\nu^2$ cannot be removed by
cyclotomic or Tate twists at any auxiliary coefficient prime other
than $3$ and $37$.  The equality case for a type induced from an
imaginary quadratic field is also unavailable.  The contemporary
completion of the Eisenstein-congruence argument in
\citet[Theorem~1.6 and Corollary~1.8, pp.~4--5]{BurungaleHeTianYe2026}
preserves this divisibility direction.

Finally, $37$ is irregular, with irregular index $32$ and
$B_{32}=-7709321041217/510$.  Thus the regular-prime Selmer
equalities in \citet[\S4.4.2]{ShuExpanded2021} do not apply.
These are failures of the specified hypotheses, not a claim that no
other arithmetic theorem could apply.

The image of the Fermat curve in $P$ is nonconstant and generates
$P$: the cyclic quotient and the projection above induce a surjection
of Jacobians.  It is therefore not a torsion translate of an elliptic
subvariety.  Proving rank zero and whole-$\Sha$ finiteness for $P$
would supply a proper finite-descent input.  The earlier split-prime
criterion needs only rank zero together with vanishing of
$T_\ell\Sha(P/\mathbb Q)$ for all but finitely many primes completely
split in a fixed finite Galois extension.  Neither conclusion follows
from the present structural calculation.
